\medskip

Dans la figure ci-dessous, ABCDEFGH est un parallélépipède rectangle tel que 

AB $= 5$, AD $= 3$ et AE $= 2$.

L'espace est muni d'un repère orthonormé d'origine A dans lequel les points B, D et E ont respectivement pour coordonnées (5~;~0~;~0), (0~;~3~;~0) et (0~;~0~;~2).

\begin{center}
\psset{unit=1cm}
\begin{pspicture}(9,5)

\psframe(0.2,0.2)(7.2,3)
\psline(7.2,0.2)(8.4,1.6)(8.4,4.4)(7.2,3)
\psline(8.4,4.4)(1.6,4.4)(0.2,3)
\psline[linestyle=dashed](0.2,0.2)(1.6,1.6)(8.4,1.6)
\psline[linestyle=dashed](1.6,1.6)(1.6,4.4)
\pspolygon[linestyle=dotted,linewidth=1.5pt](0.2,0.2)(8.4,1.6)(3,4.4)
\uput[dl](0.2,0.2){A} \uput[dr](7,0.2){B} \uput[r](8.4,1.6){C}
\uput[ur](1.6,1.6){D} \uput[l](0.2,3){E} \uput[ul](7.2,3){F}
\uput[ur](8.4,4.4){G} \uput[ul](1.6,4.4){H} \uput[u](3,4.4){M}
\end{pspicture}
\end{center}

\begin{enumerate}
\item 
	\begin{enumerate}
		\item Donner, dans le repère considéré, les coordonnées des points H et G.
		\item Donner une représentation paramétrique de la droite (GH).
	\end{enumerate}	
\item Soit M un point du segment [GH] tel que $\vect{\text{HM}} =k\vect{\text{HG}}$ avec $k$ un nombre réel de l'intervalle [0~;~1].
	\begin{enumerate}
		\item Justifier que les coordonnées de M sont $(5k~;~3~;~2)$.
		\item En déduire que $\vect{\text{AM}} \cdot \vect{\text{CM}} = 25k^2  - 25k + 4$.
		\item Déterminer les valeurs de $k$ pour lesquelles AMC est un triangle rectangle en
M.
	\end{enumerate}
\end{enumerate}	

Dans toute la suite de l'exercice, on considère que le point M a pour coordonnées (1~;~3~;~2).

On admet que le triangle AMC est rectangle en M .

On rappelle que le volume d'un tétraèdre est donné par la formule $\dfrac13 \times\text{Aire de la base}  \times h$ où $h$ est la hauteur relative à la base.

\begin{enumerate}[resume]
\item On considère le point K de coordonnées (1~;~3~;~0).
	\begin{enumerate}
		\item Déterminer une équation cartésienne du plan (ACD).
		\item Justifier que le point K est le projeté orthogonal du point M sur le plan (ACD).
		\item En déduire le volume du tétraèdre MACD.
	\end{enumerate}	
\item On note P le projeté orthogonal du point D sur le plan (AMC).

Calculer la distance DP ; en donner une valeur arrondie à $10^{-1}$.
\end{enumerate}
